\documentclass[conference]{IEEEtran}
\IEEEoverridecommandlockouts
\usepackage{cite}
\usepackage{amsmath,amssymb,amsfonts}
% \usepackage{algorithmic}   % unused; omit unless an algorithm float is added
\usepackage{graphicx}
\usepackage{textcomp}
\usepackage{xcolor}
\usepackage{booktabs}
\usepackage{float}
\usepackage{microtype}
\usepackage{url}          % required by \url in the bibliography
\usepackage{balance}
\def\BibTeX{{\rm B\kern-.05em{\sc i\kern-.025em b}\kern-.08em
    T\kern-.1667em\lower.7ex\hbox{E}\kern-.125emX}}

\begin{document}
\title{Loading-Gain Calibration Unlocks Tall RRAM Crossbars for 2-bit Floating-Point Compute-in-Memory\\
% {\footnotesize \textsuperscript{*}Technical Paper Draft for DATE Track D15}
}
% \author{\IEEEauthorblockN{First Author, Second Author, Third Author}
% \IEEEauthorblockA{\textit{Department of Electrical and Computer Engineering} \\
% \textit{University / Institution Name}\\
% City, Country \\
% \{author1, author2, author3\}@institution.edu}
% }
\maketitle

\begin{abstract}
Analog compute-in-memory (CIM) accelerators demand aggressive weight quantization, but mapping digital formats to analog constraints remains a fundamental challenge. This paper demonstrates a structural co-design match: the 2-bit block-floating-point format (FP2-E1M0) is exactly the mathematical format that a Resistive RAM (ReRAM) device can physically store. Because multi-level ReRAM reliably supports only three conductance states, it perfectly accommodates FP2's five signed levels by utilizing three magnitudes mapped to a 2T2R differential pair to supply the sign.
While this format-device alignment is ideal, it introduces a strict architectural constraint: a crossbar column produces a single current and can therefore carry only one scale factor, physically locking the block size $B$ to the tile height $M$. The FP2 specification fixes $B=32$, placing it in direct conflict with the hardware requirement of $M \ge 128$ needed to amortise analog-to-digital converter (ADC) overheads. Forcing a taller tile naively causes severe degradation, collapsing ResNet-18 accuracy on CIFAR-10 from $92.42\%$ to $10.00\%$. Using circuit-exact simulations across $87{,}168$ tiles, we demonstrate that this failure stems from a deterministic shared-sense-resistor loading gain, $I_{\text{act}} = I_{\text{ideal}}/(1 + R_s G_{\text{col}})$, rather than standard additive noise. By applying the reciprocal of this compile-time constant to each bitline before differential subtraction, our proposed digital calibration exactly cancels the attenuation. This data-free correction decouples $M$ from $B$, enabling a $4\times$ taller tile and reducing the ADC requirement to 6 bits, which simultaneously recovers the digital accuracy ceiling and delivers a $13.1\times$ improvement in energy efficiency.
\end{abstract}

\begin{IEEEkeywords}
Compute-in-Memory, Resistive RAM, Hardware-Software Co-Design, Quantization-Aware Training, Loading Calibration, Emerging Memories.
\end{IEEEkeywords}

\section{Introduction}
Deep neural networks (DNNs) have achieved state-of-the-art accuracy across edge and cloud domains, but they are severely bounded by the ``Memory Wall'' on conventional Von Neumann architectures-- i.e., the growing speed gap in the processor and memory architectures. Analog compute-in-memory (CIM) promises to remove the dominant energy and speed cost of neural network inference by performing multiply-accumulate operations directly within the memory array through Ohm's law. This eliminates the high energy and latency overhead of fetching weights. Resistive RAM (ReRAM) is a leading candidate due to its non-volatility, density, and support for multiple analog conductance states. Realising this promise requires aggressive weight quantization, as every weight must reside physically in the cells. The primary contribution of this work is demonstrating that the 2-bit block-floating-point format (FP2-E1M0) is not merely a convenient quantization choice for CIM; it is the ideal mathematical format the device can physically hold.

Programming a ReRAM filament to land repeatably on eight distinguishable conductance levels exceeds practical device physics; three to four states is the reliable ceiling~\cite{perez2021}. The FP2-E1M0 format utilizes five signed levels: $\{-1, -0.5, 0, +0.5, +1\}$. By employing a 2T2R differential cell pair to supply the sign, these five levels require exactly three conductance magnitudes. Consequently, FP2 lands exactly on what the device supports, avoiding the 256 levels required by INT8 or the loss of dynamic range inherent to binary formats.

In a typical ReRAM crossbar, input activations are converted to analog voltages and applied to the row wires. Multiplication is executed via Ohm's Law ($I = V \cdot G$), and accumulation occurs along the vertical column bitlines via Kirchhoff's Current Law ($I_{\text{col}} = \sum V_i G_i$). However, physical non-idealities severely distort these currents, threatening the viability of analog accelerators in high-precision workloads~\cite{shanbhag}.

Furthermore, this ideal format-device match forces a severe architectural constraint. A block-floating-point format shares one scale factor across $B$ consecutive weights~\cite{fp2}. Because a crossbar column produces a single summed current, it can only carry one scale factor. A tile spanning two blocks would require applying two different scales to a sum already formed by Kirchhoff's current law, destroying the information needed to separate the contributions. Consequently, the block size $B$ and the crossbar tile height $M$ are forced to be the same physical parameter. The FP2 specification fixes $B=32$, placing it in direct conflict with the hardware requirement of $M \ge 128$ needed to amortise the area and energy cost of the ADCs.

To break this collision, this paper presents three core contributions:
\begin{itemize}
    \item We demonstrate that analog column errors are dominated by a deterministic, systematic loading-gain error ($1/(1+R_s G_{\text{col}})$) rather than stochastic noise, and provide a detailed circuit-level nodal derivation.
    \item We introduce a per-column compile-time digital calibration flow that mathematically divides out this attenuation factor on dual-bitline 2T2R cells, validated against $87{,}168$ network tiles on ngspice.
    \item We co-design this calibration with Quantization-Aware Training (QAT) and Successive Approximation Register (SAR) ADC sweeps to show that a flat 6-bit ADC resolution is optimal across all column heights, achieving 81.3 TOPS/W while proving extreme robustness to device variability and conductance drift.
\end{itemize}

\section{Background}
\subsection{2T2R ReRAM Cells}
A ReRAM cell consists of a thin insulating oxide layer sandwiched between two metal electrodes. Applying a sufficient voltage forms a conductive microscopic filament across the oxide, switching the cell to a Low Resistance State (LRS), while applying a reverse voltage breaks the filament, returning it to a High Resistance State (HRS)~\cite{zahoor2020}. These states can be seen in Fig.~\ref{fig:reram_states}. By defining the input voltage $V$ as an activation and the cell's conductance $G = 1/R$ as a weight, Ohm's law ($I = V \cdot G$) natively performs multiplication in the analog domain.

Because physical conductance cannot be negative, signed weights are mapped using a differential 2T2R (two transistors, two resistors) architecture. A single weight $w_i$ is distributed across a positive cell $G^p_i$ and a negative cell $G^n_i$, such that $w_i \propto G^p_i - G^n_i$. While analog cells can theoretically hold many values, programming a filament to land repeatably on eight or more distinguishable conductance levels exceeds practical device physics. Therefore, we define three highly reliable physical states derived from compact model calibration: $R_{\text{LRS}} = 542.8\,\Omega$ ($G_{\text{LRS}} \approx 1842.3\,\mu\text{S}$), $R_{\text{MID}} = 1099.8\,\Omega$ ($G_{\text{MID}} \approx 909.3\,\mu\text{S}$), and $R_{\text{HRS}} = 218.6\,\text{k}\Omega$ ($G_{\text{HRS}} \approx 4.6\,\mu\text{S}$).

\begin{figure}[t]
    \centering
    \includegraphics[width=0.48\textwidth]{ReRAM_States.png}
    \caption{\textbf{The ReRAM resistance states.}}
    \label{fig:reram_states}
\end{figure}

\subsection{The FP2-E1M0 Data Format}
Deep learning models require aggressive quantization to fit within limited on-chip memory footprints. The FP2-E1M0 format utilizes 1 exponent bit, 0 mantissa bits, and a sign bit to yield five logarithmically spaced levels: $\ell \in \{-1, -0.5, 0, 0.5, 1\}$. To represent a wider dynamic range, this format employs block floating point, sharing an 8-bit scale factor $s$ (an E8M0 exponent) across a block of $B=32$ consecutive weights~\cite{fp2}. Mathematically, the actual weight is expressed as:
$$w_i = s \times \ell_i$$
where the shared scale constraint yields an effective storage cost of only 2.25 bits per weight.

By placing the magnitude on the bitline corresponding to the weight's polarity in the 2T2R structure, FP2's five signed levels require exactly three physical conductance magnitudes. A level of $+1.0$ is mapped as $(G^p, G^n) = (G_{\text{LRS}}, G_{\text{HRS}})$, $+0.5$ as $(G_{\text{MID}}, G_{\text{HRS}})$, and $0.0$ as $(G_{\text{HRS}}, G_{\text{HRS}})$. Negative levels invert this assignment. This creates a perfect structural match with ReRAM's reliable three-state limit.

However, mapping mathematical ideals to physical devices introduces specific non-idealities. First, $R_{\text{MID}}/R_{\text{LRS}} = 2.0262$, not exactly 2. Second, finite HRS conductance shifts every level. Consequently, the $+1.0$ level reads $0.9975$, and the $+0.5$ level lands at $0.4911$. This mapping error is correctable at compile time by either trimming $R_{\text{MID}}$ to ${\sim}1082.9\,\Omega$ or folding the measured ratio into the block scale.

\subsection{CNN Inference on Crossbar Arrays}
A crossbar of $M$ rows and $K$ columns natively computes vector-matrix multiplications in a single electrical settling time. Driven by input activation voltages $V_i$, the array produces analog currents that accumulate along the bitlines via Kirchhoff's current law. For a given column, the differential output current is:
$$I_{\text{total}} = \sum_{i=1}^M V_i (G^p_i - G^n_i)$$
which yields the ideal matrix-vector product $\mathbf{I} = (\mathbf{G}^p - \mathbf{G}^n)^\top \mathbf{V}$.

Because reprogramming a ReRAM cell costs roughly $100\,\text{pJ}$---five orders of magnitude more than the ${\sim}0.001\,\text{pJ}$ required to read it---time-multiplexing weights is prohibitive~\cite{perez2021}. Consequently, Convolutional Neural Network (CNN) inference must be executed in a fully weight-stationary dataflow, where all network weights remain permanently fixed in the crossbar array.

\section{Physical Nodal Modeling of Ohmic Attenuation}
In a standard ReRAM crossbar, input activations $x_i$ are mapped to physical voltages $V_i$ via digital-to-analog converters (DACs) on the rows. These voltages drive the rows of a 2D grid of ReRAM cells, where the programmed conductances $G_{ij}$ act as weights. The ideal current summing at the bottom of column $j$ follows Kirchhoff's Current Law (KCL):
\begin{equation}
    I_{\text{ideal}, j} = \sum_{i=1}^{M} V_i \cdot G_{ij}
\end{equation}
To convert this analog column current back to a digital representation, a small sense resistor ($R_s$) is placed at the bottom of the column bitline, developing a sense voltage $V_{\text{sense}, j} = I_{\text{actual}, j} \cdot R_s$ that is quantized by an analog-to-digital converter (ADC). However, because $R_s > 0\,\Omega$, the column line is raised slightly above electrical ground. The potential of the column line node $V_{\text{col}, j}$ is equal to $V_{\text{sense}, j}$. This back-voltage reduces the effective voltage drop across every single ReRAM cell in that column from $V_i$ to $(V_i - V_{\text{col}, j})$.

Applying KCL at the column node yields:
\begin{equation}
    \frac{V_{\text{col}, j}}{R_s} = \sum_{i=1}^{M} (V_i - V_{\text{col}, j}) \cdot G_{ij}
\end{equation}
Distributing the terms on the right-hand side:
\begin{equation}
    \frac{V_{\text{col}, j}}{R_s} = \sum_{i=1}^{M} V_i \cdot G_{ij} - V_{\text{col}, j} \sum_{i=1}^{M} G_{ij}
\end{equation}
Recognizing that $\sum_{i=1}^{M} V_i \cdot G_{ij} = I_{\text{ideal}, j}$ and defining the total programmed column conductance as $G_{\text{col}, j} = \sum_{i=1}^{M} G_{ij}$, we rewrite the relation as:
\begin{equation}
    V_{\text{col}, j} \left( \frac{1}{R_s} + G_{\text{col}, j} \right) = I_{\text{ideal}, j}
\end{equation}
Multiplying both sides by $R_s$:
\begin{equation}
    V_{\text{col}, j} (1 + R_s \cdot G_{\text{col}, j}) = R_s \cdot I_{\text{ideal}, j}
\end{equation}
Since $V_{\text{col}, j} = I_{\text{actual}, j} \cdot R_s$, we substitute and divide out $R_s$:
\begin{equation}
    I_{\text{actual}, j} = \frac{I_{\text{ideal}, j}}{1 + R_s \cdot G_{\text{col}, j}}
    \label{eq:loading_gain}
\end{equation}
Equation~(\ref{eq:loading_gain}) proves that physical column current is attenuated by a systematic, column-conductance-dependent loading-gain factor:
\begin{equation}
    \eta_j = \frac{1}{1 + R_s \cdot G_{\text{col}, j}}
\end{equation}
As physical columns are scaled taller to amortize ADC power, the sum of conductances $G_{\text{col}, j}$ scales linearly with $M$. Our physical circuit simulations track this current retention collapse, as shown in Fig.~\ref{fig:loading_attenuation}.

\begin{figure}[t]
    \centering
    \includegraphics[width=0.48\textwidth]{fig2_loading_attenuation.png}
    \caption{\textbf{Systematic Ohmic Attenuation}: Left axis plots the physical current retention fraction ($\%$) as a function of physical tile height $M$, demonstrating a drop from 73.9\% at $M=32$ to a catastrophic 26.0\% at $M=256$. Right axis demonstrates the corresponding collapse in raw uncalibrated ResNet-18 accuracy on CIFAR-10, highlighting the severe distortion of dot-product outputs.}
    \label{fig:loading_attenuation}
\end{figure}

\section{Compile-Time Per-Column Calibration Theory}
\label{sec:calibration}
Prevailing evaluation models device non-ideality as a zero-mean Gaussian perturbation drawn independently per cell and applied to the programmed conductance. Joshi \emph{et al.} state the assumption explicitly, casting read and write error into a single term $\delta G^{l}_{ij}$ and taking it to be ``Gaussian distributed, which is usually the case for analog memristive hardware''~\cite{joshi2020}; Roy \emph{et al.} likewise treat bitcell conductance variation and input DAC mismatch as zero-mean noise currents~\cite{shanbhag}. Such a model cannot represent ohmic loading, for two independent reasons. It is zero-mean, whereas loading is a systematic attenuation. And it is independent per cell, so its contribution to a column sum falls as $1/\sqrt{M}$, whereas the loading gain is common to every cell on the bitline and grows with $M$.

Because the weights are loaded onto the chip at compile-time and remain static during inference, the column conductance $G_{\text{col}, j}$ is a compile-time constant. This allows us to mathematically cancel out the attenuation in the digital domain.

To represent both positive and negative weights, we deploy a 2T2R differential cell architecture. The weight of cell $i$ in column $j$ is represented as the difference between its positive and negative programmed conductances:
\begin{equation}
    w_{ij} \propto G_{p, ij} - G_{n, ij}
\end{equation}
By routing the positive and negative differential columns to separate ADCs, we read out their attenuated currents independently:
\begin{equation}
    I_{\text{actual}, p, j} = \frac{I_{\text{ideal}, p, j}}{1 + R_s \cdot G_{\text{col}, p, j}}
\end{equation}
\begin{equation}
    I_{\text{actual}, n, j} = \frac{I_{\text{ideal}, n, j}}{1 + R_s \cdot G_{\text{col}, n, j}}
\end{equation}
Because $G_{\text{col}, p, j}$ and $G_{\text{col}, n, j}$ are compile-time constants calculated directly from our target weight matrix, we pre-compute the inverse loading-gain factors:
\begin{equation}
    \alpha_{p, j} = 1 + R_s \cdot G_{\text{col}, p, j}
\end{equation}
\begin{equation}
    \alpha_{n, j} = 1 + R_s \cdot G_{\text{col}, n, j}
\end{equation}
By deploying a digital multiplier at the output of the ADCs, we scale each digital branch by its corresponding compile-time factor prior to differential subtraction:
\begin{equation}
    \begin{aligned}
        I_{\text{calibrated}, j} &= \alpha_{p, j} \cdot I_{\text{actual}, p, j} - \alpha_{n, j} \cdot I_{\text{actual}, n, j} \\
        &= I_{\text{ideal}, p, j} - I_{\text{ideal}, n, j} \\
        &= (G_{p, j} - G_{n, j})^T V
    \end{aligned}
\end{equation}
The algebra collapses back to the ideal dot-product with mathematical exactness, reducing the physical loading error to less than $10^{-12}\%$ (bounded only by machine float precision). This converts the physical tile height $M$ from a binding circuit-level constraint into a free architectural design parameter.

\section{Quantization-Aware Training (QAT) Co-Design}
Although calibration cancels physical loading distortion, block-based low-bit quantization (such as the 2-bit FP2 representation) still introduces mathematical rounding errors (quantization noise) as block sizes scale. To prevent this, we co-design our calibration with a Quantization-Aware Training (QAT) pipeline.

During the forward pass of training, weights are quantized using a simulated FP2-E1M0 format with shared block-wise scaling. Gradients are calculated with respect to the quantized weights using a straight-through estimator, allowing the remaining parameters to adapt and compensate for the coarse resolution. As shown in Fig.~\ref{fig:qat_history}, fine-tuning the model for only 10 epochs completely recovers the accuracy loss across all block sizes.

\begin{figure}[t]
    \centering
    \includegraphics[width=0.48\textwidth]{fig6_qat_history.png}
    \caption{\textbf{QAT Fine-Tuning Convergence}: ResNet-18 accuracy trajectory on CIFAR-10 over 10 epochs of QAT. Test accuracy (blue) recovers rapidly to 92.77\%, while cell utilization (orange) and quantization residual (grey) remain stable. This proves that QAT adapts the network parameters around a fixed physical noise floor rather than attempting to shrink it.}
    \label{fig:qat_history}
\end{figure}

\section{Exhaustive SPICE-Backed Validation}
To guarantee the physical accuracy of our vectorized GPU nodal solver, we constructed a rigorous multi-level validation chain against \texttt{ngspice}~\cite{ngspice} (a golden reference analog DC solver). Rather than relying on simple behavioral noise injections, we solve the exact nodal equations of the 2D crossbar circuit.

\begin{figure}[t]
    \centering
    \includegraphics[width=0.48\textwidth]{fig5_ngspice_validation.png}
    \caption{\textbf{Physical Validation against SPICE}: High-resolution histogram of relative disagreement (\%) between our fast vectorized solver and the golden reference \texttt{ngspice} solver. Across 87,168 physical tiles representing 357 million MAC reads, the worst-case relative error is a microscopic $3.41 \times 10^{-7}\%$, confirming high physical fidelity.}
    \label{fig:ngspice_validation}
\end{figure}

Our validation sweep (Fig.~\ref{fig:ngspice_validation}) spanning 87,168 tiles confirms that our vectorized GPU solver is identical to SPICE down to numerical precision, validating all architectural sweeps.

\begin{figure}[t]
    \centering
    \includegraphics[width=0.48\textwidth]{fig7_layer_snr.png}
    \caption{\textbf{Layer-Wise Signal-to-Noise Ratio (SNR)}: Analysis of the 19 layers of ResNet-18 under uncalibrated analog execution (orange) versus our calibrated proposed execution (blue). Calibration boosts the output SNR by up to 6.3\,dB across individual convolutional layers, preventing signal degradation in deeper stages.}
    \label{fig:layer_snr}
\end{figure}

\section{Experimental Evaluation \& Co-Design Sweeps}
All results use the parameters of Table~\ref{tab:params}. Device resistances are obtained by SPICE calibration of a Verilog-A compact model; peripheral and memory constants are taken from published measurements at or near 65\,nm.

\begin{table}[H]
    \centering
    \caption{Simulation parameters.}
    \label{tab:params}
    \resizebox{\columnwidth}{!}{%
    \begin{tabular}{lll}
        \toprule
        \textbf{Parameter} & \textbf{Value} & \textbf{Source} \\
        \midrule
        Technology node & 65\,nm & --- \\
        $R_{\text{LRS}}/R_{\text{MID}}/R_{\text{HRS}}$ & 542.8\,/\,1099.8\,/\,218.6\,k$\Omega$ & compact model \\
        $R_s$, $V_{\text{read}}$ & 20\,$\Omega$, 0.1\,V & design choice \\
        ADC resolution, energy & 6\,b, 20\,fJ/conv-step & \cite{firlej2024} \\
        1T1R cell area & 40\,$F^2$ & \cite{shim2020} \\
        6T SRAM cell area & 140\,$F^2$ & \cite{shim2020} \\
        SRAM retention leakage & 120\,pW/bit & \cite{verma2008} \\
        DRAM access & 10\,pJ/bit & \cite{horowitz} \\
        INT8 MAC & 0.591\,pJ & \cite{horowitz} \\
        ReRAM write & ${\sim}100$\,pJ/cell & \cite{perez2021} \\
        Drift exponent $\nu$ (LRS, HRS) & 0.01, 0.05 & \cite{pries2022} \\
        Block size, shared scale & $B{=}32$, E8M0 & \cite{fp2} \\
        Chip-level area and energy & --- & \cite{neurosim} \\
        \bottomrule
    \end{tabular}%
    }
\end{table}

\subsection{Headline Comparison: Baseline vs. Proposed}
We evaluate our joint co-design framework against a strict conventional baseline on the CIFAR-10 ResNet-18 benchmark. This baseline represents an uncalibrated passive-readout crossbar built exactly as dictated by the FP2 specification ($B=M=32$) using standard 8-bit ADCs.

By applying our digital calibration to break the $B=M$ constraint and scaling the physical column height to $M=128$, the proposed co-design recovers \textbf{+14.77 points of accuracy} while delivering a \textbf{$13.1\times$ energy efficiency improvement}, fully accounting for ADC and DAC energy overheads.

\begin{table}[H]
    \centering
    \caption{Headline Performance Comparison on CIFAR-10 ResNet-18}
    \label{tab:headline}
    \resizebox{\columnwidth}{!}{%
    \begin{tabular}{lcc}
        \toprule
        \textbf{Metric} & \textbf{Baseline (Raw)} & \textbf{Proposed (Calibrated)} \\
        \midrule
        Block Size / Tile Height ($B=M$) & 32 & \textbf{128} \\
        Per-Column Gain Calibration & No & \textbf{Yes} \\
        ADC Resolution & 8-bit & \textbf{6-bit} \\
        \midrule
        CIFAR-10 Top-1 Accuracy & 77.66\% & \textbf{92.43\%} \\
        Gap to Ideal FP32 (92.81\%) & -15.15 pts & \textbf{-0.38 pts} \\
        Energy per MAC & 0.323 pJ & \textbf{0.0237 pJ} \\
        Energy Efficiency (Tile-Level) & 6.2 TOPS/W & \textbf{81.3 TOPS/W} \\
        \bottomrule
    \end{tabular}%
    }
\end{table}

\subsection{Optimal Sensing Resistance and SNR}
\label{sec:rsense}
Independent analytical frameworks establish that crossbar compute SNR balances ADC clipping and quantization noise at an optimal sensing resistance, identified as $R_s^* = 835\,\Omega$ for ReRAM~\cite{shanbhag}. An uncalibrated array cannot operate at this optimum; at $R_s = 835\,\Omega$, the raw readout collapses entirely to chance ($10.00\%$) due to extreme loading attenuation. Because our proposed digital calibration exactly cancels the loading term, it frees the sensing resistance constraint. At $835\,\Omega$, the calibrated array recovers to $92.59\%$ accuracy (within $0.22$ points of the digital ceiling), allowing the architecture to operate at its theoretical SNR optimum. The same framework predicts that compute SNR depends on the resistive contrast $R_{\text{off}}/R_{\text{on}}$ rather than on the absolute value of $R_{\text{on}}$, saturating once contrast exceeds 12--15; our device sweep, run with the contrast held at 403, confirms this independently, holding calibrated accuracy within $0.40$ points while $R_{\text{LRS}}$ moves over $184\times$.

\subsection{Robustness to Device Variability}
\label{sec:robustness}
ReRAM cells suffer from physical manufacturing and programming variations. We model this as a log-normal conductance variation ($\sigma_{dev}$) from 0\% to 20\% across 10 random seeds. We evaluate a ``blind'' calibration model, where digital scale factors are calculated using only the nominal target conductances (completely avoiding slow, high-voltage on-chip write-verify loops).

\begin{table}[H]
    \centering
    \caption{CIFAR-10 Accuracy (\%) vs. Device Variability ($\sigma_{dev}$)}
    \label{tab:variability}
    \resizebox{\columnwidth}{!}{%
    \begin{tabular}{ccccc}
        \toprule
        \textbf{Tile Height $M$} & $\mathbf{\sigma_{dev}=0\%}$ & $\mathbf{\sigma_{dev}=10\%}$ & $\mathbf{\sigma_{dev}=20\%}$ & \textbf{Max Loss} \\
        \midrule
        32  & 92.42\% & 92.32\% & 91.88\% & -0.54 pts \\
        64  & 92.47\% & 92.34\% & 91.86\% & -0.61 pts \\
        128 & 92.81\% & 92.66\% & \textbf{92.03\%} & -0.78 pts \\
        256 & 92.43\% & 92.25\% & 91.62\% & -0.81 pts \\
        \bottomrule
    \end{tabular}%
    }
\end{table}

As shown in Table~\ref{tab:variability} and Fig.~\ref{fig:variability_robustness}, our blind calibration holds accuracy loss to \textbf{under 0.81 percentage points} everywhere, even under harsh $\sigma_{dev}=20\%$ device scatter, proving extreme physical robustness.

\begin{figure}[t]
    \centering
    \includegraphics[width=0.48\textwidth]{fig8_variability.png}
    \caption{\textbf{Robustness to Device Variability}: CIFAR-10 accuracy curves under log-normal device variation across physical tile heights $M = 32 \rightarrow 256$, proving that blind compile-time calibration remains highly stable across 10 random seeds without requiring active on-chip tuning.}
    \label{fig:variability_robustness}
\end{figure}

\subsection{Conductance Drift and Periodic Refresh}
Over time, conductive filaments relax structurally, causing conductance to decay ($G(t) = G_0(t/t_{\text{ref}})^{-\nu}$)~\cite{pries2022,joshi2020}. At $B=32$, a stale calibration computed once at manufacture drops from $92.31\%$ to $65.17\%$ after one year. However, because drift is a power law, damage accumulates with the logarithm of time. Periodically recomputing the calibration from the drifted array (requiring only $\sim\!10$ refreshes over a product lifetime) holds accuracy flat at $91.87\%$.

\subsection{ADC Resolution Exploration}
ADCs dominate both the energy and the area budget of an analog CIM tile, and their share is exactly what tile height buys back (Table~\ref{tab:adc_share}). Because SAR conversion energy scales as $2^{N}$, right-sizing the resolution is the second lever on the same term.

\begin{table}[H]
    \centering
    \caption{Measured ADC share of tile energy and tile area against tile height. Energy shares are MAC-weighted across the 19 convolutional layers; the per-layer spread is given in brackets.}
    \label{tab:adc_share}
    \begin{tabular}{crrr}
        \toprule
        $B{=}M$ & \textbf{ADC energy} & \textbf{ADC area} & \textbf{pJ/MAC} \\
        \midrule
        32  & 90.3\% \;[83.9--93.5] & 55.2\% & 0.0887 \\
        64  & 87.2\% \;[81.2--90.2] & 40.6\% & 0.0459 \\
        128 & \textbf{81.6\%} \;[76.4--84.3] & \textbf{26.5\%} & \textbf{0.0246} \\
        256 & 72.4\% \;[68.2--74.5] & 15.7\% & 0.0141 \\
        \bottomrule
    \end{tabular}
\end{table}

The converters are 90.3\% of tile energy at $B{=}32$ and 81.6\% at $B{=}128$; their area share more than halves over the same range, from 55.2\% to 26.5\%. Energy per MAC falls $3.6\times$ from $B{=}32$ to $B{=}128$. This is the amortisation that the $B{=}M$ collision would otherwise forbid. We swept the ADC resolution from 4 to 10 bits across all column heights.

\begin{table}[H]
    \centering
    \caption{Corrected CIFAR-10 Accuracy vs. ADC Resolution}
    \label{tab:adc_sweep}
    \begin{tabular}{cccccc}
        \toprule
        \textbf{Tile Height $M$} & \textbf{4-bit} & \textbf{5-bit} & \textbf{6-bit} & \textbf{8-bit} & \textbf{Exact} \\
        \midrule
        32  & 74.88\% & 91.27\% & \textbf{92.41\%} & 92.47\% & 92.42\% \\
        64  & 66.93\% & 90.97\% & \textbf{92.27\%} & 92.47\% & 92.47\% \\
        128 & 77.18\% & 91.43\% & \textbf{92.43\%} & 92.70\% & 92.81\% \\
        256 & 77.88\% & 90.85\% & \textbf{92.01\%} & 92.49\% & 92.43\% \\
        \bottomrule
    \end{tabular}
\end{table}

The sweep (Table~\ref{tab:adc_sweep} and Fig.~\ref{fig:adc_sweep}) proves that accuracy saturates flatly at \textbf{6 bits} across all tile heights. Dropping from a conventional 8-bit converter to our optimized 6-bit configuration cuts ADC energy by four times. Converter energy is priced at 20\,fJ/conversion-step, bracketed by measured 65\,nm SAR converters operating at 40--80\,MSps~\cite{firlej2024}.

\begin{figure}[t]
    \centering
    \includegraphics[width=0.48\textwidth]{fig10_adc_resolution.png}
    \caption{\textbf{ADC Resolution Sweep}: ResNet-18 accuracy as a function of ADC bit depth across varying block sizes, proving that accuracy saturates flatly at 6 bits for all column heights, preventing amortized energy from being handed back.}
    \label{fig:adc_sweep}
\end{figure}

\begin{figure}[t]
    \centering
    \includegraphics[width=0.48\textwidth]{fig3_energy_breakdown.png}
    \caption{\textbf{System-Level Energy Breakdown}: Stacked energy allocation (pJ/MAC) for uncalibrated and calibrated architectures, demonstrating how scaling column heights amortizes ADC energy (grey block) and DAC energy (orange block) to boost efficiency.}
    \label{fig:energy_breakdown}
\end{figure}

\subsection{Memory Footprint and Density Benchmarks}
We compared our FP2-ReRAM design to static RAM (SRAM) technologies for a model storing 11.16M weights. SRAM cells require 6 transistors per bit, resulting in a massive physical footprint and continuous leakage.

\begin{table}[H]
    \centering
    \caption{Memory Technology Active Area and Standby Leakage}
    \label{tab:memory_bench}
    \begin{tabular}{lccc}
        \toprule
        \textbf{Technology} & \textbf{Data Format} & \textbf{Active Area} & \textbf{Standby Leakage} \\
        \midrule
        Digital SRAM & FP32 & 211.20 mm$^2$ & $>40$ mW \\
        Compressed SRAM & FP2 & 13.20 mm$^2$ & 2.68 mW \\
        \textbf{Proposed ReRAM} & \textbf{FP2} & \textbf{3.77 mm$^2$} & \textbf{0.00 mW} \\
        \bottomrule
    \end{tabular}
\end{table}

At the weight array, the proposed FP2-ReRAM design is $3.5\times$ denser than FP2 in SRAM and $56\times$ denser than FP32 in SRAM (Fig.~\ref{fig:memory_area}), and draws no standby leakage to hold weights indefinitely. Cell areas are measured values: $40\,F^2$ for the 1T1R cell, bracketed by $31\,F^2$ at 90\,nm~\cite{shim2020} and $84\,F^2$ at 40\,nm, and $140\,F^2$ for the 6T SRAM cell against a measured range of $150$--$200\,F^2$~\cite{shim2020}, so the SRAM columns are understated and the comparison is conservative. Leakage is from a 65\,nm measurement~\cite{verma2008}. These are array-level figures: they count cells and cell area and exclude converters, interconnect and buffers, a distinction returned to in Section~\ref{sec:limits}.

\begin{figure}[t]
    \centering
    \includegraphics[width=0.48\textwidth]{fig9_memory_area.png}
    \caption{\textbf{Silicon Active Area Footprint}: Horizontal bar chart comparing active silicon area ($mm^2$) required to store 11.16 million weights across digital FP32 SRAM, compressed FP2 SRAM, and proposed dual-bitline 2T2R ReRAM, demonstrating a 56$\times$ reduction at the weight array.}
    \label{fig:memory_area}
\end{figure}

\begin{figure}[t]
    \centering
    \includegraphics[width=0.48\textwidth]{fig4_pareto.png}
    \caption{\textbf{Energy-Accuracy Pareto Frontier}: CIFAR-10 Top-1 Accuracy (\%) versus Energy Efficiency (TOPS/W). Our calibrated proposed architecture (blue line) pushes the design envelope upward, delivering high-accuracy inference at extreme energy efficiency levels compared to conventional digital SRAM baselines.}
    \label{fig:pareto}
\end{figure}

\section{Related Work}
\subsection{Silicon Compute-in-Memory Macros}
Table~\ref{tab:sota} places this work against measured macros. Two cautions govern how the efficiency column should be read. First, reported efficiencies are not on a common basis: some are \emph{array-level}, counting only the multiply and accumulate inside the macro, and some are \emph{full-network}, counting buffers, interconnect and data movement. The gap is not small. Huang \emph{et al.} report both for one piece of silicon --- 918\,TOPS/W for MAC operations against 41.1\,TOPS/W for full-network inference, a factor of $22$~\cite{facim}. Second, efficiency scales strongly with process node, and the entries below span 0.18\,\textmu m to 22\,nm.

\begin{table}[H]
    \centering
    \caption{Measured compute-in-memory macros. ``Level'' distinguishes array-level from full-network efficiency; the two are not comparable.}
    \label{tab:sota}
    \resizebox{\columnwidth}{!}{%
    \begin{tabular}{lllrl}
        \toprule
        \textbf{Work} & \textbf{Node} & \textbf{Precision} & \textbf{TOPS/W} & \textbf{Level} \\
        \midrule
        \multicolumn{5}{l}{\emph{ReRAM}} \\
        Liu, ISSCC'20~\cite{swtt}                 & 130\,nm & 3\,b W, 1\,b act & 78.4 & array \\
        Hung, JSSC'23~\cite{hung2023}             & 22\,nm  & 8\,b/8\,b/19\,b  & 21.6 & array \\
        \midrule
        \multicolumn{5}{l}{\emph{SRAM}} \\
        Huang, JETCAS'26~\cite{facim}             & 0.18\,\textmu m & INT8 & 918  & array \\
        Huang, JETCAS'26~\cite{facim}             & 0.18\,\textmu m & INT8 & 41.1 & network \\
        Challagundla, JETCAS'26~\cite{tdccim}     & 28\,nm & INT8 & 38.46 & network \\
        Lu, TVLSI'26~\cite{chargeshare}           & 28\,nm & multibit & 30.74 & array \\
        Zou, TCAS-II'26~\cite{sramlut}            & 28\,nm & INT8 dense  & 2.32  & array \\
        Zou, TCAS-II'26~\cite{sramlut}            & 28\,nm & INT8 sparse & 72.93 & array \\
        Zhang, ISCAS'26~\cite{infp}               & 45\,nm & BF16 & 17.97$^{\dagger}$ & array \\
        \midrule
        \textbf{This work} & 65\,nm & \textbf{FP2} & \textbf{81.3} & \textbf{tile} \\
        \textbf{This work} & 65\,nm & \textbf{FP2} & \textbf{47.1} & \textbf{chip} \\
        \bottomrule
        \multicolumn{5}{l}{\footnotesize $^{\dagger}$\,TFLOPS/W.} \\
    \end{tabular}%
    }
\end{table}

The nearest architectural neighbour is the sign-weighted 2T2R array of Liu \emph{et al.}~\cite{swtt}, which places positive and negative weights on the same column so that matched conductances cancel in common mode, suppressing IR drop and lowering power $1.9\times$ against a 1T1R array on the same die. That design is differential 2T2R and targets a loading-type error, as here. It differs in what it does about the residual: common-mode cancellation makes the remaining drop small, whereas the correction in Section~\ref{sec:calibration} measures it and divides it out identically. The two compose. Their demonstration is a two-layer perceptron on MNIST at 130\,nm; the network evaluated here is a 19-layer convolutional network on CIFAR-10.

Roy \emph{et al.}~\cite{shanbhag} give an analytical framework for the compute SNR of resistive crossbars, against which Section~\ref{sec:rsense} is evaluated. This work is a correction method assessed within that framework, not a competing limits analysis.

\subsection{Alternatives to Digital Calibration}
\subsubsection{Transimpedance sensing}
A transimpedance amplifier (TIA) holds the bitline at true virtual ground, so with $v_{\text{bl}}=0$ the term $1+R_sG_{\text{col}}$ never arises. The circuit fix is standard and is not disputed here. It is not free (Table~\ref{tab:tia}).

\begin{table}[H]
    \centering
    \caption{Cost of replacing passive sensing with per-bitline transimpedance amplifiers, $M{=}128$, $K{=}16$, two converters per column.}
    \label{tab:tia}
    \begin{tabular}{lrr}
        \toprule
         & \textbf{Passive $+$ calib.} & \textbf{TIA sensing} \\
        \midrule
        Extra analog blocks   & none & 32 op-amps/tile \\
        Tile area             & 75\,369\,\textmu m$^2$ & 107\,369\,\textmu m$^2$ (\textbf{+42\%}) \\
        Die area, whole net   & 411.8\,mm$^2$ & 586.6\,mm$^2$ \\
        Added static power    & 0 & 384\,\textmu W/tile (\textbf{2.10\,W}) \\
        vs.\ measured array power & --- & $\mathbf{4.5\times}$ \\
        Correction cost       & 1 multiply/column & none needed \\
        \bottomrule
    \end{tabular}
\end{table}

The static-power row decides it. 384\,\textmu W per tile is $4.5\times$ the 85.8\,\textmu W the array itself dissipates, it is always on, and across the network it adds 2.10\,W to a design whose entire compute budget is 274\,nJ per inference. The digital correction costs one multiply per column, folded into the block-scale multiply already present in the datapath.

\subsubsection{BatchNorm recalibration}
The closest retraining-free alternative recomputes BatchNorm running statistics on the degraded network, absorbing the average current deficit into the per-channel affine. It is a strong baseline: it lifts $B{=}256$ from chance to 90.24\%. It is not exact (Table~\ref{tab:bn}).

\begin{table}[H]
    \centering
    \caption{BatchNorm recalibration against per-column calibration. Full 10\,000-image test set, 6-bit ADC. Gaps are to the digital ceiling.}
    \label{tab:bn}
    \begin{tabular}{crrrrr}
        \toprule
        $B{=}M$ & \textbf{Ceiling} & \textbf{Raw} & \textbf{BN-recal} & \textbf{Ours} & \textbf{BN gap} \\
        \midrule
        32  & 92.42\% & 78.90\% & 90.75\% & \textbf{92.35\%} & $-1.67$ \\
        64  & 92.47\% & 57.47\% & 90.09\% & \textbf{92.09\%} & $-2.38$ \\
        128 & 92.81\% & 11.53\% & 90.62\% & \textbf{92.42\%} & $-2.19$ \\
        256 & 92.43\% & 10.00\% & 90.24\% & \textbf{91.88\%} & $-2.19$ \\
        \bottomrule
    \end{tabular}
\end{table}

Worst-case gap is 2.38 points for BatchNorm against 0.55 for the per-column correction, a factor of $4.3$. Three further differences are structural rather than numerical: the correction is exact rather than fitted, it requires no calibration data, and it applies to networks without BatchNorm layers at all, such as transformer projection layers. One discrepancy should be stated rather than elided: reported figures for BatchNorm recalibration under IR drop reach within 0.16\% of floating point, against 1.67--2.38 points here. The likely cause is that IR-drop attenuation is roughly uniform across a layer, whereas the shared-sense-resistor gain varies with each column's programmed conductance, leaving less for a per-channel affine to absorb.

\subsubsection{Training-time weight pre-distortion}
Non-ideality-aware training transforms the target weights into a pre-distorted matrix so that the degraded hardware output matches the ideal product~\cite{joshi2020}. This requires retraining and assumes a fixed non-ideality. The correction here is computed from the conductances \emph{as programmed}, needs no retraining, and tracks the array through variability and drift --- the regimes measured in Section~\ref{sec:robustness}.

\section{Limitations}
\label{sec:limits}
The shared-sense-resistor loading term is exactly correctable, but real interconnects contain distributed parasitics. Evaluating the array with a full 2D resistive mesh solver, wordline metal ($0.5\,\Omega$) leaves a negligible $1.82\%$ residual, whereas bitline metal ($0.5\,\Omega$) and wordline driver resistance ($50\,\Omega$) leave $49.1\%$ and $25.4\%$: distributed voltage drops are not captured by a single per-column scalar. The claim is therefore narrowed to the shared-sense-resistor term, with distributed wire parasitics remaining a physical limit.

The device parameters derive from a Verilog-A compact model rather than fabricated silicon, and every downstream number inherits that. The structure of the argument does not: the loading error is a deterministic function of programmed conductance whatever those values are. The drift exponents are the weakest input, since the power law is established for amorphous phase-change material rather than ReRAM and $\nu_{\text{LRS}}$, $\nu_{\text{HRS}}$ and $\sigma_\nu$ are assumed. We therefore sweep $\sigma_\nu$ over $0.001$--$0.016$ and report the boundary at $\approx\!0.008$ rather than a single drift cost.

Area comparisons require care with denominators. The $3.5\times$ and $56\times$ density figures are array-level: at 2-bit weights the weight array occupies under $1\%$ of a chip-level floorplan, where converters, interconnect, global buffer and pooling dominate, so substituting SRAM cells of five times the area grows the chip by $3.7\%$ rather than by $3.5\times$. The chip-level tool's total also responds to floorplan geometry rather than component sizes alone: widening the 1T1R cell four-fold with all else fixed \emph{reduces} its reported total from $120.4$ to $62.4\,\text{mm}^2$ while every named component grows or is unchanged. All areas reported here are sums of named components. Finally, a fully weight-stationary ResNet-18 at $B{=}128$ models to $412\,\text{mm}^2$ at 65\,nm, so the density argument depends on scaled nodes.

\section{Conclusion}
This paper resolved the physical-to-mathematical block size coupling in emerging ReRAM CIM crossbars. By mathematically modeling systematic ohmic bitline loading as a deterministic, compile-time constant, we cancel attenuation precisely in the digital domain. Paired with QAT and right-sized 6-bit ADCs, our architecture achieves 92.43\% accuracy on CIFAR-10 while reaching a tile-level energy efficiency of 81.3 TOPS/W on tall $M=128$ columns. This breaks the limitations of standard CIM architectures, paving the way for highly efficient, next-generation edge accelerators.

\begin{thebibliography}{99}
\balance
\bibitem{fp2}
B.~D. Rouhani \emph{et al.}, ``Microscaling data formats for deep learning,''
\emph{arXiv:2310.10537}, 2023.

\bibitem{neurosim}
X.~Peng, S.~Huang, Y.~Luo, X.~Sun, and S.~Yu, ``DNN+NeuroSim: An end-to-end
benchmarking framework for compute-in-memory accelerators with versatile device
technologies,'' in \emph{IEDM}, 2019.

\bibitem{ngspice}
ngspice, ``Open source mixed-mode circuit simulator.'' [Online]. Available:
\url{https://ngspice.sourceforge.io}

\bibitem{shanbhag}
S.~K. Roy, A.~Patil, and N.~R. Shanbhag, ``Fundamental limits on the
computational accuracy of resistive crossbar-based in-memory architectures,''
in \emph{IEEE ISCAS}, 2022, pp.\ 384--388.

\bibitem{swtt}
Q.~Liu \emph{et al.}, ``33.2 A fully integrated analog ReRAM based
78.4\,TOPS/W compute-in-memory chip with fully parallel MAC computing,'' in
\emph{IEEE ISSCC}, 2020, pp.\ 500--502.

\bibitem{horowitz}
M.~Horowitz, ``Computing's energy problem (and what we can do about it),'' in
\emph{IEEE ISSCC Dig.\ Tech.\ Papers}, 2014, pp.\ 10--14.

\bibitem{shim2020}
W.~Shim, X.~Sun, J.-S. Seo, and S.~Yu, ``Two-bit-per-cell RRAM-based in-memory
computing for area-/energy-efficient deep learning,'' \emph{IEEE Solid-State
Circuits Lett.}, vol.~3, pp.\ 288--291, 2020.

\bibitem{verma2008}
N.~Verma and A.~P. Chandrakasan, ``A 256\,kb 65\,nm 8T subthreshold SRAM
employing sense-amplifier redundancy,'' \emph{IEEE J. Solid-State Circuits},
vol.~43, no.~1, pp.\ 141--149, Jan.\ 2008.

\bibitem{perez2021}
E.~Perez \emph{et al.}, ``Variability and energy consumption tradeoffs in
multilevel programming of RRAM arrays,'' \emph{IEEE Trans. Electron Devices},
vol.~68, no.~6, pp.\ 2693--2698, Jun.\ 2021.

\bibitem{firlej2024}
M.~Firlej \emph{et al.}, ``Ultra-low power 10-bit 50--90\,MSps SAR ADCs in
65\,nm CMOS for multi-channel ASICs,'' \emph{J. Instrum.}, vol.~19, P01029,
2024.

\bibitem{facim}
W.-C. Huang, H.-T. Chien, S.-J. Chang, J.-H. Hu, and C.-H. Kuo, ``A fully
analog computing-in-memory macro with INT8-MAC operations for edge-AI
devices,'' \emph{IEEE J. Emerg. Sel. Topics Circuits Syst.}, vol.~16, no.~2,
pp.\ 239--250, Jun.\ 2026.

\bibitem{pries2022}
J.~Pries \emph{et al.}, ``Resistance drift, convergence and inversion in
amorphous phase change materials,'' \emph{Adv.\ Funct.\ Mater.}, 2022.

\bibitem{joshi2020}
V.~Joshi \emph{et al.}, ``Accurate deep neural network inference using
computational phase-change memory,'' \emph{Nat.\ Commun.}, vol.~11, 2473,
2020.

\bibitem{zahoor2020}
F.~Zahoor, T.~Z. Azni Zulkifli, and F.~A. Khanday, ``Resistive random access
memory (RRAM): an overview of materials, switching mechanism, performance,
multilevel cell (MLC) storage, modeling, and applications,'' \emph{Nanoscale
Res.\ Lett.}, vol.~15, 90, 2020.

\bibitem{hung2023}
J.-M. Hung \emph{et al.}, ``An 8-b precision 8-Mb ReRAM compute-in-memory macro
using direct-current-free time-domain readout scheme for AI edge devices,''
\emph{IEEE J. Solid-State Circuits}, vol.~58, no.~1, pp.\ 303--315, Jan.\ 2023.

\bibitem{tdccim}
D.~Challagundla, I.~Bezzam, and R.~Islam, ``TDC-CiM: Time-to-digital
converter-based resonant compute-in-memory for INT8 CNNs with layer-optimized
SRAM mapping,'' \emph{IEEE J. Emerg. Sel. Topics Circuits Syst.}, vol.~16,
no.~2, pp.\ 227--238, Jun.\ 2026.

\bibitem{chargeshare}
W.~Lu, L.~Tan, S.~Wang, T.~Xue, C.~Peng, Z.~Lin, and X.~Wu, ``A 28-nm 9-kb
SRAM computing-in-memory macro with segmented charge sharing for multimode MAC
operations,'' \emph{IEEE Trans. Very Large Scale Integr. (VLSI) Syst.},
vol.~34, no.~5, May 2026.

\bibitem{sramlut}
X.~Zou \emph{et al.}, ``A high-SNR SRAM-LUT-based in-memory-computing macro
supporting energy-efficient INT8 and FP8 computation,'' \emph{IEEE Trans.
Circuits Syst. II, Exp. Briefs}, vol.~73, no.~7, Jul.\ 2026.

\bibitem{infp}
J.~Zhang, T.~Dong, S.~Li, Z.~Yan, and S.~Huang, ``InFP: A 17.97 TFLOPS/W
reconfigurable SRAM based computing-in-memory macro for BF16 MAC operations,''
in \emph{IEEE Int. Symp. Circuits Syst. (ISCAS)}, 2026.
\end{thebibliography}

\end{document}
